\section{Effective context as controlled state}\label{sec:C11:implications}
Long-context systems expose an important asymmetry: storing more evidence is straightforward to specify, whereas identifying what will matter later requires assumptions about future tasks. A larger window relaxes a resource constraint, but it does not remove the need for organization, authority, or temporal interpretation. At sufficiently long horizons, unfiltered history becomes a record of abandoned plans, obsolete facts, and adversarial content as well as useful evidence. More history can increase both opportunity and ambiguity.

The original synthesis of this chapter is that effective context should be evaluated as a controlled state representation. Its quality depends jointly on information preserved, the policy's ability to use that information, and the operational cost of making it available. Equation~\eqref{eq:C11:informationloss} supplies a precise idealization for predictive loss, while the control-state condition extends the concern to future action consequences. Neither equation turns present agents into certified sufficient-state estimators; both clarify which claims an experiment would need to support.

Several practical rules follow from that formulation. Preserve authoritative constraints independently of conversational salience. Retain references to raw evidence whenever summaries may be revised. Expose version and timestamp metadata at the decision point instead of expecting the model to infer them from prose. Measure complete-task cost rather than one-call context size. Treat cache state as an implementation of inference, not as a validated semantic memory. These are design deductions under explicit state and cost assumptions rather than reported superiority of a particular platform.

The limits of current evidence are substantial. Controlled context benchmarks are useful, but their target distributions need not match a changing enterprise workflow or a long scientific investigation. Public papers can establish gains under specific environments and budgets without establishing reliable indefinite operation. A context-management method may transfer across several tasks yet still depend on the same backbone, retriever, or evaluator. Independent reproduction and operational deployment answer additional questions.

For implementation, a useful boundary is to keep immutable interaction evidence, mutable task state, and the assembled model input distinct. This allows context construction to be inspected as a transformation rather than an invisible consequence of token overflow. It also permits counterfactual replay: reconstruct a prior decision with the original context, an edited context, or an oracle evidence set, while holding the environment snapshot fixed. Such replay is expensive, but it can reveal failure mechanisms that aggregate success rates conceal.

Chapter~\ref{ch:C12} develops the editing policy that chooses the model's working view. Chapter~\ref{ch:C13} treats the event and session substrate that makes those choices recoverable. The resulting picture is an agent whose finite context is one deliberately managed component of a larger stateful system. Its generality remains an empirical question about task transfer, adaptation, and reliable execution under bounded resources.
